\section{Number Systems}
The most natural and familiar number system is the set of \textit{natural numbers}, 
\[\mathbb{N} = \{1, 2, 3, \dots\}.
\]
However, solving equations using only natural numbers is often challenging. To address this, we extend $\mathbb{N}$ by introducing $0$ and additive inverses, forming the set of \textit{integers}, 
\[
\mathbb{Z} = \{\dots, -3, -2, -1, 0, 1, 2, 3, \dots\}.
\]
$\Z$ allows for systematic techniques in algebra, particularly in solving linear equations. In $\mathbb{Z}$, we can add, subtract, and multiply, but division is not always possible. To allow division, we extend $\mathbb{Z}$ to the set of \textit{rational numbers}, 
\[
\mathbb{Q} = \left\{\frac{p}{q} \mid p, q \in \mathbb{Z}, q \neq 0, gcd(p,q)=1 \right\}.
\]
However, $\mathbb{Q}$ is still not sufficient to solve all equations. For example, $\sqrt{2}$ is not a rational number but a solution to the equation $x^2 = 2$. To include such numbers, we extend $\mathbb{Q}$ to the set of \textit{real numbers}, $\mathbb{R}$. 
Yet even $\mathbb{R}$ is not enough for certain equations, such as $x^2 + 1 = 0$ have no solutions in $\mathbb{R}$. To resolve this, we introduce the \textit{complex numbers}, 
\[
\mathbb{C} = \{a + bi \mid a, b \in \mathbb{R}, \ i^2 = -1\}.
\]
According to the \textit{Fundamental Theorem of Algebra}, $\mathbb{C}$ is \textit{algebraically closed}, meaning every polynomial equation with complex coefficients has a solution in $\mathbb{C}$. Thus, $\mathbb{C}$ is the most comprehensive number system we need for solving polynomial equations, and we do not extend it further in this context. Notice,
\[
\N \subset \Z \subset \Q \subset \R \subset \C.
\]
\section{Well-Ordering Principle}
In Peano Arithmetic, the \textit{Principle of Mathematical Induction (PMI)} is assumed as an axiom. And from this axiom, we can derive a fundamental result known as the \textit{Well-Ordering Principle}.

\textbf{Well-Ordering Principle.} Every non-empty subset of $\N$ has a least element.
\section{Goals and Motivation for doing NT}
To be written once I figure out myself. Hopefully, at the end of this course I will be able to write something here. :3